% ============================================================
%  Poul Martin Møller, "On the Concept of Irony"
%  English translation of "Om Begrebet Ironie"
%  (Efterladte Skrifter, bd. 3, Brudstykke II; printed pp. 152--158)
%
%  Mirrors transcription.tex 1:1. \opage{N} marks the point at
%  which printed page N begins. Danish quotes „…" -> ``…''.
%  Book titles -> \emph{}. Only the treatment of MORAL irony was
%  completed by the author; the piece breaks off as it turns to
%  poetic irony.
% ============================================================

\documentclass[12pt]{article}
\usepackage[utf8]{inputenc}
\usepackage[T1]{fontenc}
\usepackage{libertinus}
\usepackage{libertinust1math}
\usepackage{textalpha}
\usepackage[english]{babel}
\usepackage[protrusion=true,expansion=true]{microtype}
\usepackage[margin=1.4in]{geometry}
\usepackage{setspace}
\usepackage{fancyhdr}
\usepackage{marginnote}
\usepackage[colorlinks=true,linkcolor=black,urlcolor=black,
  pdftitle={On the Concept of Irony},
  pdfauthor={Poul Martin Møller}]{hyperref}
\setlength{\headheight}{15pt}
\pagestyle{fancy}
\fancyhf{}
\fancyhead[LE,RO]{\thepage}
\fancyhead[RE]{\textit{On the Concept of Irony}}
\fancyhead[LO]{\textit{Poul Martin Møller}}
\setlength{\parindent}{1.5em}
\setlength{\parskip}{0pt}
\onehalfspacing

\newcommand{\opage}[1]{\marginnote{\footnotesize\textit{[#1]}}}
\newcommand{\ombreak}{\par\medskip\begin{center}\rule{0.32\textwidth}{0.4pt}\end{center}\medskip}

\title{\textbf{On the Concept of Irony}}
\author{Poul Martin Møller}
\date{Posthumous Writings, vol.\ 3 (Fragment II)}

\begin{document}
\maketitle

\begin{center}\rule{0.32\textwidth}{0.4pt}\end{center}
\medskip

\opage{152}On the occasion of the observations which the esteemed author of the foregoing review has communicated concerning the concept of irony\footnote{What is meant is Provost Tryde's review of Sibbern's \emph{Æsthetik} in the \emph{Maanedsskrift f.\ Lit.}, vol.\ 13 (1835), and especially p.\ 200 ff. On making himself acquainted, as a member of the Maanedsskrift's editorial board, with Provost Tryde's manuscript, P.\ M.\ was moved to append his own remarks; but he brought only the Introduction to completion. Since it may be regarded as a whole in itself (on moral irony), it may deserve a place among his writings, all the more as it furnishes a contribution to the history of the concept in question in our literature; yet it must be remembered that the manner of treating the subject was conditioned by the requirements of an Introduction.} --- as it has taken shape in the aesthetic usage of the present day --- let it be permitted to one of the members of the editorial board to append a few aphoristic remarks, which may perhaps give some readers a clearer insight into the origin of this concept and its various transformations.

It will be excused that the beginning is fetched from somewhat far off; we shall nonetheless soon return to the matter. In the moral self-consciousness of the Greeks, as we come to know it in its first germ among the people's oldest poets --- e.g.\ in several passages of Homer --- and as we later find it in full development in Aristotle, the natural desire is regarded as the positive side of the good; reason, on the other hand, is the limiting principle, which sanctions only a certain quantum of this \opage{153}natural content and rejects another. Virtue is a measure, a middle way between two extremes. We thus find at this standpoint two different determinations in the concept of the good, which are reconciled with one another because their opposition has not been developed with real clarity. But it had necessarily to be brought to consciousness thereafter that mere impulse, as a blind, incomprehensible and irrational ground of determination, could not be sanctioned by reason. This Aristotle already, the most perfect representative of that standpoint, must have felt; and from this one may well explain the peculiarity that, according to his conceptions of knowledge, the content of practical philosophy cannot become an object of science in the higher sense of the word.

Of the two heterogeneous and supposedly irreconcilable elements in the concept of the good, now the one, now the other has by turns been championed with one-sidedness in the history of science. Now the attempt to determine the aim of life by rational thought has been wholly given up, and it has been held to be wisdom to let the natural desire in its accidental shape be the governing power; now, by a thinking that kept solely to the individual standpoint, men have sought to derive the moral maxims they found before them from a dogmatic proposition or from a wholly abstract and empty rule. Thus one has seen the wholly contradictory precepts: ``Let us return to the state of nature'' (Rousseau), and: ``One must forsake the natural constitution'' (Hobbes) --- made to prevail one after the other.

The endeavour to determine the content of the good solely from the standpoint of pure reason, or by a merely a priori thinking, has been carried through with the greatest consistency by the leading men of British philosophy. In this philosophy men have attempted, from merely formal principles, to derive the circle of ethical maxims prevailing in society. All heteronomy, every merely natural ground of determination, they wished to reject, and by pure \opage{154}a priori thinking to specify the content of moral self-consciousness. But notwithstanding the enthusiasm with which the thought of absolute autonomy in such a sense filled the adherents of this school, such a plan could nevertheless not possibly be carried through scientifically. Both Kant and Fichte manifestly see themselves compelled to make use of heteronomous maxims in order to procure a real content for their abstract principles.

Now, when it was taught that the natural ends were to be carried out as means for the moral law, but not for the sake of the pleasure that unavoidably attached itself to their execution, there had, by the nature of the case, to be sought a method by which it could be decided which end was, in every single moment, the right means whereby the universal law could be put into effect. Here Fichte now placed the highest authority in the consciousness of the individual, so that this, according to its own moral conviction in the particular case, was to decide what was its duty. This standpoint, which must always retain its (relative) validity, thus ascribes to the individual the capacity to recognise immediately according to which of the various rules the universal is to be realised in the particular case. ---

Therefore, from this moral idealism, when it is maintained in a one-sided way and not united with the recognition that rational commands are already realised outside the individual, objectionable consequences may be drawn. Subjective conviction then comes to be regarded as the highest, in that the will of the individual is identified with the moral law. The individual will is then set above every particular law and above the dead letter of every determinate rule, which has validity only when it is sanctioned by the individual will. This will, autonomous and exalted above everything, is then consistently regarded as raised above the civil laws, which have validity only where they agree with it --- a view that permits all crimes that may be committed, and have often been com\opage{155}mitted, with a good conscience and youthful, moral enthusiasm. Thus there is nothing which, in itself and independently of the moral subject, is good; but this [subject], as identical with the moral law, can sanctify any action whatever by its conviction; for that which otherwise, in itself, should be evil ceases to be so when the will, which is the main thing, is noble, or when the intention was good.

From this theory the transition is easily made to another, which sets the subjective will above the moral law; for it is indeed, practically, really set above it, when that which is in itself evil can be sanctified by the will. The moral law has at bottom lost its sanctity, and the individual makes, if not ethical, yet at least logical progress by expressly assuming itself to be raised above the moral law. This theory is found expressly stated in F.\ Schlegel's youthful work \emph{Lucinde}, as impudent as it is brilliant, which is one of the most luxuriant offshoots of that inconsiderate, audacious self-confidence which is sufficient unto itself. We cite, in confirmation of what has been said, a couple of utterances from \emph{Lucinde} itself: ``Nothing is more absurd than when moralists now reproach one another for egoism! for what god can be venerable to the man who is not his own god?'' Further, that ``in the enjoyment of existence one must raise oneself above all finite ends and purposes''; that the highest is therefore the divine idleness, and that only in the holy vacancy of genuine passivity can one recall one's whole self --- which pure vegetating is then the true life and the true religion. --- This standpoint Hegel has taken up in his practical philosophy and designated by the name irony, which he explains as the ``subjectivity that knows itself as the highest.'' Herein he recognises nothing positive but the activity of uttering this consciousness for others, and that is plainly a renunciation of all endeavours in the objective world.

Although this practical resignation is infinitely different from \opage{156}the energetic morality taught in Fichte's works (which it is doubtless superfluous to remark), it is nevertheless certain that it descends from Fichtean idealism, as it is also known that Fred.\ Schlegel was one of the most enthusiastic admirers of the Fichtean doctrine of science. In particular, irony is a consistent continuation of the fruitless striving to bring a self-enclosed moral system into being solely from the standpoint of the individual. This procedure had necessarily to end in the lack of all content, in a moral nihilism.

The irony here spoken of thus designates not a merely conceivable standpoint in the development of moral consciousness; nor has it come forth as an inexplicable notion upon which some single eccentric man was accidentally brought, but it was a necessary result of an erring train of thought, which showed itself in a similar shape among several contemporaries.

Even among us this aberration, so deeply grounded in the development of the age, has had its fainter echo, although, owing to the modesty and caution of our people, it has left few or no traces in our literature. But in conversation one has often been able to hear that view resound fairly closely; and even if it has not been expressed with full assurance and self-confidence, it has nevertheless been uttered as a result to which one would be led if one wished to abandon oneself to the course of a prejudice-free thinking. It has very commonly shown itself as a kind of distrust toward the concept of morality, which one preferred to avoid in scientific discourse, since one no longer found in it an adequate expression for practical freedom; one has also now and then in poetry made sport with the concepts moral, morality, and moral, and expressed a dim consciousness that things did not quite hang together with them. Indeed, one or another has also made use of the audacious turn that true ability would not let itself be restricted by the narrow moral law.

A shape of moral consciousness closely re\opage{157}lated to irony is morbid sentimentality: it can at least lead to the same result, namely to the opinion that man can be raised above his earthly duties. An idle longing for a better life then brings the sentimental man to give up all the definite ends of life, because everything here below is only perishable and affords no worthy stage for one who has a sense for something higher. At this standpoint a separation is made between a mere finite on the one side, which has no share whatever in infinity, and, separated from it by a yawning abyss, a wholly indeterminate higher life, which has been called the Better or the Eternal. Since now the former is the sentimental man's only worthy aim, he thinks he must, in passive resignation, with the surrender of all definite endeavour, await the coming of the Higher to him. One may here apply what Fichte, as far as is recalled, has somewhere said: Many hold that one need only let oneself be buried in order to attain a blessed life. --- This sentimentality resembles irony in this, that both consider themselves raised above all commands of duty: the latter, because he sets his individual arbitrariness as the highest; the former, because he regards all finite action in time and space as a lower sphere of activity and believes that he defiles the purity of his inner life by occupying himself with its contact. Both thus hold themselves too good to have duties: the ironist, because his individual I is the Absolute; the sentimental man, because his sphere of activity in life is not outside him; he lives only an inner life, which is filled with longings and premonitions of the future, which is the only real. But they differ from one another in this, that the ironist regards his subjectivity as the highest, while the sentimental man takes the Absolute as a One outside his person, as outside the whole visible world. Yet this difference is somewhat evened out by the fact that the sentimental man's consciousness is the seat of a recognition of the nothingness of the finite, of the reality of the indeterminate Higher, whereby he \opage{158}comes to stand far above all who with zeal pursue definite practical ends.

The ironist too can, according to his system, if he will, without being hindered by any moral straitjacket, abandon himself to the coarsest enjoyment of sensuality --- a consequence recognised in \emph{Lucinde}. This properly contradicts the sentimental theory; but since the doctrine here cannot possibly be followed out in its whole consistency, and action cannot be wholly given up so long as life lasts, and consequently contact with finitude cannot be wholly avoided, the sentimental man can, more easily than one who reveres a rational ethics, justify to himself preferring the most sensual kind among the actions that are in themselves indifferent. (When it is not worth the trouble to do the good, neither is it worth the trouble to refrain from the evil.)

But this was a digression, which was merely to serve to make clear the practical irony hitherto spoken of by bringing it together with the other species of practical nihilism, with which it agrees in several respects. One of the many meanings in which the word irony has been taken in more recent philosophy seems hereby to be sufficiently explained. Thus Hegel uses the word irony, and thereby designates one of ``the moral forms of evil.'' Whether one now considers the word well or ill chosen, its object deserves the attention of science, and the place which it occupies in the history of consciousness has been indisputably rightly assigned to it by Hegel. We now pass over to poetic irony. --- ---

\ombreak

\end{document}
